\par\Needspace{6\baselineskip}\medskip
\SourceLeaf{213}\hypertarget{native-213-article157}{}\pdfbookmark[2]{157.}{article-157}
\begin{center}157.\end{center}
Primo nunc situm circulorum maximorum a locis heliocentricis terrae ad locos geocentricos planetae ductorum determinabimus. Intersectionibus horum circulorum cum aequatore, aut si mavis illorum nodis ascendentibus, characteres $\mathfrak A,\mathfrak A',\mathfrak A''$ adscriptos concipimus, distantiasque punctorum $B,B',B''$ ab his punctis per $\Delta,\Delta',\Delta''$ designamus. In maiori operationum parte pro $A,A',A''$ iam $\mathfrak A,\mathfrak A',\mathfrak A''$, et pro $\delta,\delta',\delta''$ iam $\Delta,\Delta',\Delta''$ substituere oportebit; ubi vero $A,A',A'',\delta,\delta',\delta''$ retinere oporteat, lector attentus vel nobis non monentibus facile intelliget.

Calculo facto iam invenimus
\begin{center}
\begingroup\footnotesize\setlength{\tabcolsep}{3pt}\renewcommand{\arraystretch}{1.25}
\begin{tabular}{l|r|r|r}
\shortstack{Ascens. recta punctorum\\$\mathfrak A,\mathfrak A',\mathfrak A''$} & $233^0\,54'\,57''{,}10$ & $253^0\,8'\,57''{,}01$ & $276^0\,40'\,25''{,}87$\\
$\gamma,\gamma',\gamma''$ & $51\quad17\quad15{,}74$ & $90\quad1\quad3{,}19$ & $131\quad59\quad58{,}03$\\
$\Delta,\Delta',\Delta''$ & $215\quad58\quad49{,}27$ & $212\quad52\quad48{,}96$ & $220\quad9\quad12{,}96$\\
$\delta,\delta',\delta''$ & $56\quad26\quad34{,}19$ & $55\quad26\quad31{,}79$ & $69\quad10\quad57{,}84$\\
$\mathfrak A'D,\mathfrak AD',\mathfrak AD''$ & $23\quad54\quad52{,}13$ & $30\quad18\quad3{,}25$ & $29\quad8\quad43{,}32$\\
$\mathfrak A''D,\mathfrak A''D',\mathfrak A'D''$ & $33\quad3\quad26{,}35$ & $31\quad59\quad21{,}14$ & $22\quad20\quad6{,}91$\\
$\epsilon,\epsilon',\epsilon''$ & $47\quad1\quad54{,}69$ & $89\quad34\quad57{,}07$ & $42\quad33\quad41{,}17$\\
logarithmi sinuum & $9{,}864\,3525$ & $9{,}999\,9885$ & $9{,}830\,1910$\\
$\log\sin\tfrac12\epsilon'$ & & $9{,}847\,8971$ &\\
$\log\cos\tfrac12\epsilon'$ & & $9{,}851\,0614$ &
\end{tabular}\endgroup
\end{center}

In calculo art. 138 pro $l'$ ascensio recta puncti $\mathfrak A'$ adhibebitur. Sic invenitur
\begin{center}\begin{tabular}{ll}
$\log T\sin t$ .\,.\,.\,. & $8{,}486\,8236_n$\\
$\log T\cos t$ .\,.\,. & $9{,}284\,8162_n$.
\end{tabular}\end{center}
\noindent Hinc $t=189^0\,2'\,48''{,}83$, $\log T=9{,}290\,2527$; porro $t+\gamma'=279^0\,3'\,52''{,}02$,
\begin{center}\begin{tabular}{ll}
$\log S$ .\,.\,.\,.\,.\,.\,.\,.\,.\,.\,. & $9{,}011\,0566_n$\\
$\log T\sin(t+\gamma')$ .\,.\,. & $9{,}284\,7950_n$,
\end{tabular}\end{center}
\noindent unde $\Delta'-\sigma=208^0\,1'\,55''{,}64$, atque $\sigma=4^0\,50'\,53''{,}32$.

In formulis art. 140 pro $a$, $b$ et $\dfrac b a$ ipsos $\sin\delta$, $\sin\delta'$, $\sin\delta''$ retinere oportet, et perinde in formulis art. 142. Ad hos calculos habemus
\SourceLeaf{214}
\begin{center}\begingroup\footnotesize\setlength{\tabcolsep}{4pt}\renewcommand{\arraystretch}{1.3}
\begin{tabular}{ll}
$\mathfrak A''D'-\Delta''=171^0\,50'\,8''{,}18;$ & \shortstack{$\log\sin$ .\,.\,. $9{,}152\,3306;$\\$\log\cos$ .\,. $9{,}995\,5759_n$}\\
$\mathfrak AD'-\Delta=174\quad19\quad13{,}98;$ & \shortstack{--- .\,.\,. $8{,}995\,4722;$\\--- .\,.\,. $9{,}997\,8629_n$}\\
$\mathfrak A''D-\Delta''=172\quad54\quad13{,}39;$ & --- .\,.\,. $9{,}091\,7972$\\
$\mathfrak A'D-\Delta'+\sigma=175\quad52\quad56{,}49;$ & --- .\,.\,. $8{,}856\,1520$\\
$\mathfrak AD''-\Delta=173\quad9\quad54{,}05;$ & --- .\,.\,. $9{,}075\,5844$\\
$\mathfrak A'D''-\Delta'+\sigma=174\quad18\quad11{,}27;$ & --- .\,.\,. $8{,}996\,7978$.
\end{tabular}\endgroup\end{center}
\noindent Hinc elicimus
\[
\begin{aligned}
\log\varkappa&=0{,}921\,1850&\log\lambda&=0{,}081\,2057_n\\
\log\varkappa''&=0{,}811\,2762&\log\lambda''&=0{,}031\,9691_n\\
\log a&=0{,}109\,9088,&a&=+1{,}287\,9790\\
\log b&=0{,}181\,0404&&\\
\log\frac b a&=0{,}071\,1314&&
\end{aligned}
\]
\noindent unde fit $\log b=0{,}181\,0402$. Inter hos duos valores tantum non aequales medium $\log b=0{,}181\,0403$ adoptabimus. Denique prodit
\[
\begin{aligned}
\log c&=1{,}045\,0295\\
d&=+0{,}448\,9906\\
\log e&=9{,}210\,2894,
\end{aligned}
\]
\noindent quo pacto calculi praeliminares absoluti sunt.

Temporis intervallum inter observationem secundam et tertiam est dierum $39{,}874\,409$, inter primam et secundam dierum $30{,}900\,961$; hinc fit $\log\theta=9{,}836\,2757$, $\log\theta''=9{,}725\,5533$. Statuimus itaque \textit{ad hypothesin primam}
\[
\begin{aligned}
x=\log P&=9{,}889\,2776\\
y=\log Q&=9{,}561\,8290.
\end{aligned}
\]
\noindent Praecipua dein calculi momenta haec prodeunt:
\[
\begin{aligned}
\omega+\sigma&=20^0\,8'\,46''{,}72\\
\log Qc\sin\omega&=0{,}028\,2028.
\end{aligned}
\]
\noindent Hinc fit valor verus ipsius $z=21^0\,11'\,24''{,}30$, atque $\log r'=0{,}350\,9379$. Tres reliqui valores ipsius $z$ aequationi IV art. 141 satisfacientes in hoc casu fiunt
\[
\begin{aligned}
z&=63^0\,40'\,50''\\
z&=101\quad12\quad58\\
z&=199\quad24\quad7,
\end{aligned}
\]
\noindent\SourceLeaf{215}e quibus primus tamquam approximatio ad orbitam terrestrem spectandus est, cuius quidem aberratio, propter nimium temporis intervallum, longe hic maior est, quam in exemplo praecedente. --- E calculo ulteriori sequentes numeri resultant:
\begin{center}\begingroup\footnotesize\setlength{\tabcolsep}{6pt}\renewcommand{\arraystretch}{1.3}
\begin{tabular}{ll}
$\zeta=195^0\,12'\,2''{,}48$ & $\zeta''=196^0\,57'\,50''{,}78$\\
$\log r=0{,}364\,7022$ & $\log r''=0{,}335\,5758$\\
$\tfrac12(u''+u)=266^0\,47'\,50''{,}47$ & $\tfrac12(u''-u)=-43^0\,39'\,5''{,}33$
\end{tabular}\endgroup\end{center}
\[
\begin{aligned}
2f'&=22^0\,32'\,40''{,}86\\
2f&=13\quad5\quad41{,}17\\
2f''&=9\quad27\quad0{,}05.
\end{aligned}
\]
\noindent Differentiam inter $2f'$ et $2f+2f''$, quae hic est $0''{,}36$, inter $2f$ et $2f''$ ita dispertiemur, ut statuamus $2f=13^0\,5'\,40''{,}96$, $2f''=9^0\,26'\,59''{,}90$.

Corrigenda iam sunt tempora propter aberrationem, ubi in formulis art. 145 statuendum est $AD'-\zeta=\mathfrak AD'-\Delta+\delta-\zeta$, $A''D'-\zeta''=\mathfrak A''D'-\Delta''+\delta''-\zeta''$. Habemus itaque
\begin{center}\begin{tabular}{ll}
$\log r$ .\,.\,.\,.\,.\,.\,.\,.\,.\,.\,.\,. & $0{,}36470$\\
$\log\sin(AD'-\zeta)$ .\,.\,. & $9{,}76462$\\
$\mathrm C.\log\sin\delta$ .\,.\,.\,.\,.\,.\,. & $0{,}07918$\\
$\log\mathrm{const.}$ .\,.\,.\,.\,.\,.\,.\,. & $7{,}75633$\\\hline
& $7{,}96483$\\
reductio temporis$=$ & $0{,}009222$
\end{tabular}\end{center}
\begin{center}\begin{tabular}{ll}
$\log r'$ .\,.\,.\,.\,.\,. & $0{,}35094$\\
$\log\sin(\delta'-z)$ & $9{,}75038$\\
$\mathrm C.\log\sin\delta'$ . & $0{,}08431$\\
& $7{,}75633$\\\hline
& $7{,}94196$\\
& $0{,}008749$
\end{tabular}\end{center}
\begin{center}\begin{tabular}{ll}
$\log r''$ .\,.\,.\,.\,.\,.\,.\,.\,.\,. & $0{,}33558$\\
$\log\sin(A''D'-\zeta'')$ & $9{,}84220$\\
$\mathrm C.\log\sin\delta''$ .\,.\,.\,.\,. & $0{,}02932$\\
& $7{,}75633$\\\hline
& $7{,}96343$\\
& $0{,}009192$
\end{tabular}\end{center}
\noindent Hinc prodeunt
\begin{center}\begin{tabular}{r|r|r}
Tempora correcta & intervalla & logarithmi\\
Nov.\ $5{,}564\,852$ & &\\
& $30{,}901\,434$ & $1{,}489\,9785$\\
$36{,}466\,286$ & &\\
& $39{,}873\,966$ & $1{,}60\overset{.}{0}\,6894$\\
$76{,}340\,252$ & &
\end{tabular}\end{center}
\noindent unde derivantur logarithmi correcti quantitatum $\theta,\theta''$ resp. $9{,}836\,2708$ atque $9{,}725\,5599$. Incipiendo dein calculum elementorum ex $r',r'',2f,\theta$, prodit $\log\eta=0{,}003\,1921$, sicuti ex $r,r',2f'',\theta''$ obtinemus $\log\eta''=0{,}001\,7300$. Hinc colligitur $\log P'=9{,}890\,7512$, $\log Q'=9{,}571\,2864$, adeoque
\[X=+0{,}001\,4736,\qquad Y=+0{,}009\,4574.\]

Praecipua momenta \textit{hypothesis secundae}, in qua statuimus
\SourceLeaf{216}
\[
\begin{aligned}
x=\log P&=9{,}890\,7512\\
y=\log Q&=9{,}571\,2864,
\end{aligned}
\]
\noindent haec sunt:
\begin{center}\begingroup\footnotesize\setlength{\tabcolsep}{6pt}\renewcommand{\arraystretch}{1.3}
\begin{tabular}{r|r}
$\omega+\sigma=20^0\,8'\,0''{,}87$ & $\log r''=0{,}336\,9708$\\
$\log Qc\sin\omega=0{,}037\,3071$ & $\tfrac12(u''+u)=267^0\,6'\,10''{,}75$\\
$z=21^0\,12'\,6''{,}09$ & $\tfrac12(u''-u)=-43\quad39\quad4{,}00$\\
$\log r'=0{,}350\,7110$ & $2f'=22\quad32\quad8{,}69$\\
$\zeta=195^0\,16'\,59''{,}90$ & $2f=13\quad1\quad54{,}65$\\
$\zeta''=196\quad52\quad40{,}63$ & $2f''=9\quad30\quad14{,}38$\\
$\log r=0{,}363\,0642$ &
\end{tabular}\endgroup\end{center}
\noindent Differentia $0''{,}34$ inter $2f'$ et $2f+2f''$ ita distribuenda est, ut statuatur $2f=13^0\,1'\,54''{,}45$, $2f''=9^0\,30'\,14''{,}24$.

Si operae pretium videtur correctiones temporum hic denuo computare, invenietur pro observatione prima $0{,}009\,169$, pro secunda $0{,}008\,742$, pro tertia $0{,}009\,236$, adeoque tempora correcta Nov.\ $5{,}564\,905$, Nov.\ $36{,}466\,293$, Nov.\ $76{,}340\,208$. Hinc fit
\[
\begin{aligned}
\log\theta&=9{,}836\,2703,&\log\theta''&=9{,}725\,5594\\
\log\eta&=0{,}003\,1790,&\log\eta''&=0{,}001\,7413\\
\log P'&=9{,}890\,7268,&\log Q'&=9{,}571\,0593.
\end{aligned}
\]
\noindent Hoc itaque modo ex hypothesi secunda resultat
\[X=-0{,}000\,0244,\qquad Y=-0{,}000\,2271.\]

Denique in \textit{hypothesi tertia}, in qua statuimus
\[
\begin{aligned}
x=\log P&=9{,}890\,7268\\
y=\log Q&=9{,}571\,0593,
\end{aligned}
\]
\noindent praecipua calculi momenta ita se habent:
\begin{center}\begingroup\footnotesize\setlength{\tabcolsep}{6pt}\renewcommand{\arraystretch}{1.3}
\begin{tabular}{r|r}
$\omega+\sigma=20^0\,8'\,1''{,}62$ & $\log r''=0{,}336\,9536$\\
$\log Qc\sin\omega=0{,}037\,0857$ & $\tfrac12(u''+u)=267^0\,5'\,53''{,}09$\\
$z=21^0\,12'\,4''{,}60$ & $\tfrac12(u''-u)=-43\quad39\quad4{,}19$\\
$\log r'=0{,}350\,7191$ & $2f'=22\quad32\quad7{,}67$\\
$\zeta=195^0\,16'\,54''{,}08$ & $2f=13\quad1\quad57{,}42$\\
$\zeta''=196\quad52\quad44{,}45$ & $2f''=9\quad30\quad10{,}55$\\
$\log r=0{,}363\,0960$ &
\end{tabular}\endgroup\end{center}
\noindent\SourceLeaf{217}Differentia $0''{,}30$ hic ita distribuetur, ut statuatur $2f=13^0\,1'\,57''{,}25$, $2f''=9^0\,30'\,10''{,}42$\footnote{Haecce differentia maiuscula, in omnibusque hypothesibus tantum non aequalis, ad maximam partem inde orta est, quod $\sigma$ duabus \underline{f}ere partibus centesimis minuti secundi iusto minor, logarithmusque ipsius $b$ aliquot unitatibus iusto maior erutus erat.}.

Quum differentiae omnium horum numerorum ab iis, quos hypothesis secunda suppeditaverat, levissimae sint, tuto iam concludere licebit, hypothesin tertiam nulla amplius correctione opus habituram, adeoque hypothesin novam superfluam esse. Quocirca nunc ad calculum elementorum ex $2f',\theta',r,r''$ progredi licebit: qui quum operationibus supra amplissime iam explicatis contineatur, elementa ipsa inde resultantia in eorum gratiam, qui proprio marte eum exsequi cupient, hic apposuisse sufficiet:
\begin{center}\begingroup\footnotesize\setlength{\tabcolsep}{3pt}\renewcommand{\arraystretch}{1.3}
\begin{tabular}{ll}
\shortstack{Ascensio recta nodi ascendentis in aequatore\\.\,.\,.\,.\,.\,.\,.\,.\,.\,.\,.\,.\,.\,.\,.\,.} & $158^0\,40'\,38''{,}93$\\
\shortstack{Inclinatio orbitae ad aequatorem\\.\,.\,.\,.\,.\,.\,.\,.\,.\,.\,.\,.\,.\,.\,.\,.\,.\,.\,.\,.\,.\,.\,.\,.\,.\,.\,.} & $11\quad42\quad49{,}13$\\
\shortstack{Distantia perihelii a nodo illo ascendente\\.\,.\,.\,.\,.\,.\,.\,.\,.\,.\,.\,.\,.\,.\,.\,.\,.\,.\,.} & $323\quad14\quad56{,}92$\\
\shortstack{Anomalia media pro epocha 1806\\.\,.\,.\,.\,.\,.\,.\,.\,.\,.\,.\,.\,.\,.\,.\,.\,.\,.\,.\,.\,.\,.\,.\,.\,.} & $335\quad4\quad13{,}05$\\
\shortstack{Motus medius (sidereus) diurnus\\.\,.\,.\,.\,.\,.\,.\,.\,.\,.\,.\,.\,.\,.\,.\,.\,.\,.\,.\,.\,.\,.\,.\,.\,.\,.\,.\,.\,.\,.\,.} & $770''{,}2663$\\
\shortstack{Angulus $\varphi$\\.\,.\,.\,.\,.\,.\,.\,.\,.\,.\,.\,.\,.\,.\,.\,.\,.\,.\,.\,.\,.\,.\,.\,.\,.\,.\,.\,.\,.\,.\,.\,.\,.\,.\,.\,.\,.\,.\,.\,.\,.\,.\,.\,.\,.\,.} & $14^0\,9'\,3''{,}91$\\
\shortstack{Logarithmus semiaxis maioris\\.\,.\,.\,.\,.\,.\,.\,.\,.\,.\,.\,.\,.\,.\,.\,.\,.\,.\,.\,.\,.\,.\,.\,.\,.\,.\,.\,.\,.\,.\,.} & $0{,}442\,2438$.
\end{tabular}\endgroup\end{center}
